\BeleskeID{08-s038}
% element: 08-s038:e03
% element: 08-s038:notes
Vremena uspona i pada na izlazu u principu se razlikuju od odgovarajućih trajanja na ulazu. Samo merenje izlazne ivice ne izdvaja doprinos kola: konačno trajanje ulazne ivice utiče i na izmereno kašnjenje.

Prikazane aproksimativne formule razdvajaju izmereno kašnjenje od doprinosa kola pri idealnoj ulaznoj ivici, označenog indeksom 0. Za invertor se u kašnjenju silazne izlazne ivice pojavljuje ulazno vreme uspona $t_r$, a u kašnjenju uzlazne izlazne ivice ulazno vreme pada $t_f$. To su aproksimacije u okviru opisanog modela, a ne nova definicija kataloškog kašnjenja.
